%% ma: L11-B16-0092
%% lop: 11
%% chuong: 5
%% bai: 16
%% dang: 11.16.4
%% loai: TL
%% muc: VD
%% dap_an: "-"
%% nguon: "Ảnh giáo viên gửi 25/11/2026: Giới hạn hàm số dạng 0/0 (3 bài, có ĐS), Bài 2 câu 10"
%% kien_thuc: "Lớp 11 Bài 16 (giới hạn của hàm số tại một điểm); phân tích đa thức thành nhân tử, hằng đẳng thức"
%% trang_thai: chinh-thuc
%% ghi_chu: "ĐS gốc 1/9 sai, đúng là 2/9."
\begin{ex}
Tính giới hạn $A=\lim\limits_{x\to1}\left(\dfrac1{x^2+x-2}-\dfrac1{x^3-1}\right)$.
\loigiai{$\dfrac{x^2+x+1-(x+2)}{(x-1)(x+2)(x^2+x+1)}=\dfrac{(x-1)(x+1)}{(x-1)(x+2)(x^2+x+1)}=\dfrac{x+1}{(x+2)(x^2+x+1)}$. Vậy $A=\dfrac2{3\cdot3}=\dfrac29$.}
\end{ex}
